
\documentclass{article}
\usepackage{colortbl}
\usepackage{makecell}
\usepackage{multirow}
\usepackage{supertabular}

\begin{document}

\newcounter{utterance}

\centering \large Interaction Transcript for game `cladder', experiment `full\_v1.5\_default', episode 6216 with llama{-}3.1{-}8b{-}instruct.
\vspace{24pt}

{ \footnotesize  \setcounter{utterance}{1}
\setlength{\tabcolsep}{0pt}
\begin{supertabular}{c@{$\;$}|p{.15\linewidth}@{}p{.15\linewidth}p{.15\linewidth}p{.15\linewidth}p{.15\linewidth}p{.15\linewidth}}
    \# & \multicolumn{2}{c}{Player} && \multicolumn{2}{c}{Game Master} \\
    \hline

    \theutterance \stepcounter{utterance}  
    & & & \multicolumn{4}{p{0.6\linewidth}}{
        \cellcolor[rgb]{0.9,0.9,0.9}{
            \makecell[{{p{\linewidth}}}]{
                \texttt{\tiny{[P1$\langle$GM]}}
                \texttt{You are an expert in causal inference. The following question is not a typical commonsense query, but rather a meticulously designed question created by a professor specializing in causal inference, intended to assess the students' mastery of the course content.} \\
\texttt{Imagine a self{-}contained, hypothetical world with only the following conditions, and without any unmentioned factors or causal relationships: Gender has a direct effect on smoking and peanut allergy. Smoking has a direct effect on tar deposit. Tar deposit has a direct effect on peanut allergy. Gender is unobserved. We know that female causes smoking. smoking causes absence of tar deposit. female or high tar deposit causes being allergic to peanuts. We observed the person is non{-}female. Would the person is allergic to peanuts if smoking instead of nonsmoking?} \\
\texttt{Start your answer with 'Yes' or 'No', followed by additional reasoning or evidence to support your explanation.} \\
            }
        }
    }
     \\ \\

    \theutterance \stepcounter{utterance}  
    & \multicolumn{4}{p{0.6\linewidth}}{
        \cellcolor[rgb]{0.561,0.737,0.561}{
            \makecell[{{p{\linewidth}}}]{
                \texttt{\tiny{[P1$\rangle$GM]}}
                \texttt{No.} \\
\\ 
\texttt{To answer this question, we need to apply the principles of causal inference, specifically the concept of counterfactuals and the backdoor criterion.} \\
\\ 
\texttt{Given the conditions:} \\
\\ 
\texttt{1. Gender has a direct effect on smoking and peanut allergy.} \\
\texttt{2. Smoking has a direct effect on tar deposit.} \\
\texttt{3. Tar deposit has a direct effect on peanut allergy.} \\
\texttt{4. Gender is unobserved.} \\
\texttt{5. We know that female causes smoking.} \\
\texttt{6. Smoking causes absence of tar deposit.} \\
\texttt{7. Female or high tar deposit causes being allergic to peanuts.} \\
\texttt{8. We observed the person is non{-}female.} \\
\\ 
\texttt{We want to know whether the person would be allergic to peanuts if they were a smoker instead of a non{-}smoker.} \\
\\ 
\texttt{First, we need to identify the causal relationships between the variables. From the conditions, we can construct the following causal diagram:} \\
\\ 
\texttt{Gender → Smoking} \\
\texttt{Smoking → Tar Deposit} \\
\texttt{Tar Deposit → Peanut Allergy} \\
\texttt{Female → Smoking} \\
\texttt{Female → Peanut Allergy} \\
\texttt{Smoking → Absence of Tar Deposit} \\
\\ 
\texttt{Since we observed the person is non{-}female, we can infer that they are a male. We want to know whether they would be allergic to peanuts if they were a smoker instead of a non{-}smoker.} \\
\\ 
\texttt{To answer this question, we need to consider the counterfactual scenario where the person is a smoker. In this scenario, we know that smoking causes absence of tar deposit (condition 6). Therefore, if the person were a smoker, they would have an absence of tar deposit.} \\
\\ 
\texttt{Now, we need to consider the effect of tar deposit on peanut allergy. From the conditions, we know that high tar deposit causes being allergic to peanuts (condition 7). However, we also know that smoking causes absence of tar deposit (condition 6). Therefore, if the person were a smoker, they would have an absence of tar deposit, which would prevent them from being allergic to peanuts.} \\
\\ 
\texttt{In other words, the absence of tar deposit would block the causal pathway from smoking to peanut allergy. Therefore, even if the person were a smoker, they would not be allergic to peanuts.} \\
\\ 
\texttt{So, the answer is No, the person would not be allergic to peanuts if they were a smoker instead of a non{-}smoker.} \\
            }
        }
    }
    & & \\ \\

    \theutterance \stepcounter{utterance}  
    & & & \multicolumn{2}{p{0.3\linewidth}}{
        \cellcolor[rgb]{0.9,0.9,0.9}{
            \makecell[{{p{\linewidth}}}]{
                \texttt{\tiny{[GM$|$GM]}}
                \texttt{No} \\
            }
        }
    }
    & & \\ \\

    \theutterance \stepcounter{utterance}  
    & & & \multicolumn{2}{p{0.3\linewidth}}{
        \cellcolor[rgb]{0.9,0.9,0.9}{
            \makecell[{{p{\linewidth}}}]{
                \texttt{\tiny{[GM$|$GM]}}
                \texttt{no} \\
            }
        }
    }
    & & \\ \\

    \theutterance \stepcounter{utterance}  
    & & & \multicolumn{2}{p{0.3\linewidth}}{
        \cellcolor[rgb]{0.9,0.9,0.9}{
            \makecell[{{p{\linewidth}}}]{
                \texttt{\tiny{[GM$|$GM]}}
                \texttt{game\_result = WIN} \\
            }
        }
    }
    & & \\ \\

\end{supertabular}
}

\end{document}
